\chapter{Golden Dataset e Regress\~ao de Qualidade}
\label{chap:golden_dataset_regressao}

\chapterepigraph{``Odours have a power of persuasion stronger than that of words.''}{Patrick S\"uskind, \textit{Perfume: The Story of a Murderer} (1985).}
{\small\itshape\color{AccentGray}
S\"uskind escreveu sobre fragr\^ancias, n\~ao sobre dados. Mas a intui\c{c}\~ao \'{e} a mesma: h\'{a} sinais que convencem sem precisar ser traduzidos em palavras. Um \textit{golden dataset} funciona assim --- a cole\c{c}\~ao certa de casos de teste tem um poder de persu\textit{as}\~ao sobre a qualidade do sistema que nenhuma m\'{e}trica isolada consegue reproduzir. Quando o pipeline passa nessa bateria, a confian\c{c}a n\~ao \'{e} inferida por c\'{a}lculo; ela \'{e} demonstrada por evid\^encia.\par}
\vspace{0.6em}
% ---------------------------------------------------------------
\section{Pr\'e-requisitos}
% ---------------------------------------------------------------

\begin{itemize}[leftmargin=*]
  \item Cap\'itulo~\ref{chap:avaliacao_metricas_mercado} ---
        ROUGE-L e thresholds de qualidade;
  \item Cap\'itulo~\ref{chap:llm_as_judge} --- LLM Judge e rubrica;
  \item Conceito de integra\c{c}\~ao cont\'inua (CI/CD).
\end{itemize}

% ---------------------------------------------------------------
\section{Escopo do cap\'itulo}
% ---------------------------------------------------------------

\begin{quote}\itshape
\textbf{Onde estamos na hist\'oria:} loop externo, regress\~ao
autom\'atica. As m\'etricas e o judge est\~ao prontos. Falta o
elo final: um conjunto curado de casos de teste que detecta
regress\~oes antes que cheguem a produ\c{c}\~ao. O golden dataset
\'e a mem\'oria coletiva de qualidade do sistema.
\end{quote}

Este cap\'itulo cobre: (1) o que \'e um golden dataset e como
constru\'i-lo; (2) execu\c{c}\~ao de regress\~ao com
\texttt{golden\_dataset.py}; (3) interpreta\c{c}\~ao do diff
(regress\~oes vs.\@ melhorias); (4) integra\c{c}\~ao com CI/CD.

% ---------------------------------------------------------------
\section{Objetivos de aprendizagem}
% ---------------------------------------------------------------

Ao concluir este cap\'itulo, o leitor ser\'a capaz de:

\begin{itemize}[leftmargin=*]
  \item Definir os crit\'erios de curadoria de um golden dataset;
  \item Criar casos com thresholds por pergunta;
  \item Executar uma su\'ite de regress\~ao e interpretar o
        relat\'orio;
  \item Identificar a diferen\c{c}a entre regress\~ao e melhoria
        no diff de qualidade.
\end{itemize}

% ---------------------------------------------------------------
\section{Conceitos te\'oricos}
% ---------------------------------------------------------------

\subsection{O que \'e um golden dataset?}

Um \textit{golden dataset}\index{golden dataset|textbf} \'e um
conjunto de casos de teste curados manualmente, com:

\begin{itemize}[leftmargin=*]
  \item \textbf{Pergunta:} representativa de uso real em produ\c{c}\~ao;
  \item \textbf{Resposta de refer\^encia:} aprovada por especialistas
        de dom\'inio (analistas de produto ou compliance);
  \item \textbf{Thresholds individuais:} cada caso tem seu pr\'oprio
        limiar ROUGE-L e score m\'inimo do judge, refletindo o
        n\'ivel de fidelidade esperado para aquele tipo de pergunta.
\end{itemize}

Ao contr\'ario de benchmarks p\'ublicos, o golden dataset \'e
\textbf{propriet\'ario e confidencial}: cont\'em a expectativa de
qualidade negociada entre engenharia e neg\'ocio.

\subsection{Curadoria: sintético vs.\@ humano}

\begin{table}[htbp]
\centering
\caption{Compara\c{c}\~ao de estrat\'egias de curadoria.}
\label{tab:curadoria}
\begin{tabular}{lll}
\toprule
\textbf{Estrat\'egia} & \textbf{Vantagem} & \textbf{Limita\c{c}\~ao} \\
\midrule
Manual (humano)  & alta fidelidade de dom\'inio  & caro, lento \\
Sint\'etico (LLM) & escal\'avel, r\'apido         & pode herdar vi\'eses do LLM \\
H\'ibrido        & equil\'ibrio de custo/qualidade & requer revis\~ao humana \\
\bottomrule
\end{tabular}
\end{table}

O projeto adota a estrat\'egia h\'ibrida: perguntas geradas
sinteticamente com respostas de refer\^encia revisadas por
especialistas.

\subsection{Regress\~ao vs.\@ melhoria}

O diff compara os resultados atuais com os resultados do baseline
salvo. Um caso pode transitar de tr\^es formas:

\begin{itemize}[leftmargin=*]
  \item \textbf{Regress\~ao:} passou antes, falhou agora
        $\to$ bloquear release;
  \item \textbf{Melhoria:} falhou antes, passou agora
        $\to$ atualizar threshold;
  \item \textbf{Estabilidade:} mesmo resultado $\to$ sem a\c{c}\~ao.
\end{itemize}

% ---------------------------------------------------------------
\section{Implementa\c{c}\~ao orientada por passos}
% ---------------------------------------------------------------

\subsection{Passo 1 --- Estrutura de um caso do golden dataset}

\begin{lstlisting}[language=Python,
  caption={Estrutura de um caso no \texttt{golden\_dataset.py}.},
  label={lst:golden_caso}]
CASOS = [
    {
        "id": "GD-001",
        "segmento": "PF",
        "pergunta": "Qual e a rentabilidade do CDB prefixado de 12 meses?",
        "referencia": (
            "O CDB prefixado de 12 meses oferece rentabilidade de 12,5% ao ano, "
            "garantida pelo FGC ate R$ 250.000."
        ),
        "threshold_rouge": 0.35,
        "threshold_judge": 7,
    },
    # ... mais casos
]
\end{lstlisting}

\subsection{Passo 2 --- Criar e salvar o dataset}

\begin{lstlisting}[language=bash,
  caption={Cria\c{c}\~ao do golden dataset com baseline salvo.},
  label={lst:golden_criar}]
python golden_dataset.py --modo criar --dataset golden_dataset.jsonl
python golden_dataset.py --modo rodar --dataset golden_dataset.jsonl --saida golden_baseline.json

# Saida esperada:
# Golden dataset criado com 4 casos.
# Resultados salvos em golden_baseline.json
\end{lstlisting}

\subsection{Passo 3 --- Executar regress\~ao}

\begin{lstlisting}[language=bash,
  caption={Execu\c{c}\~ao da su\'ite de regress\~ao.},
  label={lst:golden_rodar}]
python golden_dataset.py --modo rodar --dataset golden_dataset.jsonl --saida resultados_golden.json

# Saida esperada:
# GD-001: rouge_l=0.2857  judge=10  STATUS: FALHOU (rouge abaixo de 0.35)
# GD-002: rouge_l=0.5000  judge=7   STATUS: PASSOU
# GD-003: rouge_l=0.4000  judge=7   STATUS: PASSOU
# GD-004: rouge_l=0.2500  judge=10  STATUS: FALHOU (rouge abaixo de 0.30)
# Resultado: 2/4 casos aprovados
\end{lstlisting}

O comportamento acima \'e \textbf{intencional}: GD-001 e GD-004
t\^em respostas mock propositalmente incompletas, validando que
o sistema de thresholds funciona corretamente.

\subsection{Passo 4 --- Interpretar o diff}

\begin{lstlisting}[language=bash,
  caption={Comparar execu\c{c}\~ao atual com baseline.},
  label={lst:golden_diff}]
python golden_dataset.py --modo diff --baseline golden_baseline.json --atual resultados_golden.json

# Saida esperada quando nenhuma mudanca ocorreu:
# "regressoes": []
# "melhorias": []
# "estavel": ["GD-001", "GD-002", "GD-003", "GD-004"]
\end{lstlisting}

\subsection{Passo 5 --- Integra\c{c}\~ao com CI/CD}

A su\'ite de regress\~ao deve ser executada automaticamente
em qualquer pull request que modifique prompts, modelos ou
pipeline RAG:

\begin{lstlisting}[language=bash,
  caption={Exemplo de gate de qualidade em GitHub Actions.},
  label={lst:golden_cicd}]
- name: Regressao de qualidade
  run: |
    python golden_dataset.py --modo rodar
    if python golden_dataset.py --modo diff | grep -q "Regressoes: [^0]"; then
      echo "FALHA: regressao detectada -- bloquear release"
      exit 1
    fi
\end{lstlisting}

% ---------------------------------------------------------------
\section{Autoavalia\c{c}\~ao}
% ---------------------------------------------------------------

\begin{enumerate}[leftmargin=*]
  \item Um caso que passou no baseline e falhou agora \'e
        classificado como o qu\^e?\\
        \textbf{Gabarito:} \textbf{regress\~ao} --- deve bloquear
        o release at\'e a causa ser identificada.

  \item Por que cada caso pode ter um threshold ROUGE-L
        diferente?\\
        \textbf{Gabarito:} porque o n\'ivel de fidelidade lexical
        esperado varia com o tipo de pergunta --- perguntas factuais
        (taxas, prazos) t\^em threshold mais alto; perguntas
        explicativas aceitam mais parafrasea\c{c}\~ao.

  \item \'E seguro mover um caso do golden dataset para um
        threshold mais baixo ap\'os uma regress\~ao?\\
        \textbf{Gabarito: N\~ao} --- isso mascararia a degrada\c{c}\~ao;
        o correto \'e investigar e corrigir a causa.
\end{enumerate}

% ---------------------------------------------------------------
\section{Resumo}
% ---------------------------------------------------------------

\begin{itemize}[leftmargin=*]
  \item O golden dataset \'e a mem\'oria coletiva de qualidade:
        expectativas negociadas entre engenharia e neg\'ocio.
  \item Cada caso tem thresholds individuais para ROUGE-L e
        LLM Judge.
  \item O diff classifica mudan\c{c}as em regress\~oes, melhorias
        e casos est\'aveis.
  \item Regress\~oes devem bloquear o release automaticamente no
        CI/CD.
  \item A estrat\'egia h\'ibrida (sint\'etico + revis\~ao humana)
        equilibra velocidade e fidelidade de dom\'inio.
\end{itemize}

% ---------------------------------------------------------------
\section{Plano de testes do cap\'itulo}
% ---------------------------------------------------------------

\begin{itemize}[leftmargin=*]
  \item \textbf{TC-14-01:} modo \texttt{criar} deve gerar
        \texttt{golden\_baseline.json} com 4 chaves.
  \item \textbf{TC-14-02:} modo \texttt{rodar} deve retornar
        2/4 aprovados com as respostas mock padr\~ao.
  \item \textbf{TC-14-03:} modo \texttt{diff} com baseline
        id\^entico deve reportar 0 regress\~oes.
  \item \textbf{TC-14-04:} caso com resposta vazia deve ser
        reprovado independentemente do threshold.
\end{itemize}

% ---------------------------------------------------------------
\section{Crit\'erios de aceite}
% ---------------------------------------------------------------

\begin{itemize}[leftmargin=*]
  \item TC-14-01 a TC-14-03 passam com \texttt{golden\_dataset.py}.
  \item Nenhuma exce\c{c}\~ao nos tr\^es modos sem credenciais AWS.
\end{itemize}

% ---------------------------------------------------------------
\section{Espa\c{c}o para expans\~ao iterativa}
% ---------------------------------------------------------------

\begin{itemize}[leftmargin=*]
  \item \textbf{v14.1:} ampliar para 50 casos por segmento
        (PF/PJ/FP) usando gera\c{c}\~ao sint\'etica com Claude.
  \item \textbf{v14.2:} dashboard de evolu\c{c}\~ao hist\'orica
        de scores ao longo de deploys.
  \item \textbf{Lacuna:} protocolo de expira\c{c}\~ao de casos ---
        como desativar casos que n\~ao refletem mais pol\'iticas
        de produto atuais.
\end{itemize}

\section{Objetivos de aprendizagem}
\begin{itemize}[leftmargin=*]
  \item Entender os conceitos nucleares do tema.
  \item Aplicar o tema em uma implementacao concreta do projeto.
  \item Definir evidencias objetivas de qualidade e funcionamento.
\end{itemize}

\section{Implementacao orientada por passos}
\begin{enumerate}[leftmargin=*]
  \item Preparacao do ambiente e insumos.
  \item Implementacao guiada do fluxo principal.
  \item Validacao dos resultados esperados.
  \item Registro de decisoes tecnicas e trade-offs.
\end{enumerate}

\section{Plano de testes do capitulo}
\subsection{Teste funcional}
\begin{itemize}[leftmargin=*]
  \item Definir casos de entrada e saida esperada.
  \item Verificar comportamento nominal e casos de erro.
\end{itemize}

\subsection{Teste de desempenho}
\begin{itemize}[leftmargin=*]
  \item Medir latencia ponta a ponta.
  \item Medir uso estimado de tokens por execucao.
\end{itemize}

\subsection{Teste de qualidade de resposta}
\begin{itemize}[leftmargin=*]
  \item Medir aderencia ao contexto recuperado.
  \item Medir consistencia da resposta com o objetivo do usuario.
\end{itemize}

\section{Criterios de aceite}
\begin{itemize}[leftmargin=*]
  \item Implementacao executa sem erros bloqueantes.
  \item Testes definidos no capitulo atingem resultado esperado.
  \item Evidencias de execucao sao anexadas pelos autores.
\end{itemize}

\section{Espaco para expansao iterativa}
\begin{itemize}[leftmargin=*]
  \item Notas de versao do capitulo.
  \item Melhorias planejadas.
  \item Lacunas para investigacao futura.
\end{itemize}
